\section{SayCan：让 LLM 计划，让 Affordance 决定能不能做}

RT-1 扩展了 skill library，却没有回答“给我清理桌面”怎样被拆成多步技能。SayCan 把语言模型当作 high-level planner，同时用机器人 value / affordance model 约束每一步是否在当前状态可执行，从而连接 Internet-scale common sense 与 grounded control。

\subsection{从 Skill Learning 转向 Planning}

Agenda 的第三项进入 SayCan。课堂 timeline 随后把 2021--2022 的 multi-task robot policies 与 2022--2023 的 LLM-powered planning 连接起来：低层动作仍来自 robot data，高层组合开始借用语言模型。

\lecturefigure{slide-23-agenda-saycan.jpg}{Agenda：SayCan 进入 language-conditioned planning 层}{官方录像恢复幻灯片 V023；Stanford CS25 Lecture 15}

\lecturefigure{slide-24-robotics-foundation-timeline.jpg}{机器人研究 timeline：从 scale real operations 到用 LLM 加速机器人系统}{官方录像恢复幻灯片 V024；Stanford CS25 Lecture 15}

\subsection{两个问题：Fixed Commands 与 No Grounding}

第一，机器人只能调用有限技能集合，例如 `find apple`、`pick apple`、`go table`；LLM 必须学会“说机器人语言”。第二，LLM 没有摄像头和动力学，不知道苹果是否存在、夹爪是否能抓到；系统必须给语言先验加 physical affordance grounding。

\lecturefigure{slide-25-language-models-robotics-challenges.jpg}{Language models for robotics：固定 skill vocabulary 与缺少 real-world grounding}{官方录像恢复幻灯片 V025；Stanford CS25 Lecture 15}

\begin{warningbox}{Planner 不能创造不存在的 Skill}
LLM 可以把目标分解得很漂亮，但若 low-level library 没有“开抽屉”或“擦拭液体”，planner 只能输出不可执行文本。系统上限由 grounded skills 的覆盖决定，这也是讲者在 Q\&A 中明确指出的 SayCan bottleneck。
\end{warningbox}

\subsection{SayCan Score：Useful 与 Possible 的乘积}

对候选技能 $k$、高层目标 $g$ 和当前状态 $s$，SayCan 可写成
\begin{equation}
\operatorname{Score}(k\mid g,s)
\propto p_{\mathrm{LM}}(k\mid g,h)
\cdot p_{\mathrm{aff}}(\mathrm{success}\mid s,k),
\end{equation}
其中 $h$ 是已执行步骤。第一项问“这一步对完成语言目标是否有用”，第二项问“机器人现在是否做得到”。在 log space 中两者相加，避免某个语言上合理但物理上不可能的动作占据首位。

\lecturefigure{slide-26-saycan-architecture.jpg}{SayCan：LLM usefulness 与 robotic affordance 联合选择下一项技能}{官方录像恢复幻灯片 V026；Ahn et al., 2022}

\begin{knowledgebox}{读图：候选排序而不是自由生成控制}
中间表列出 `find apple`、`pick apple` 等技能。LLM 分数偏好符合目标的步骤，affordance 分数过滤当前不可执行项，系统选择 combined score 最高者，执行后重新排序。低层 controller 不由 LLM token 直接驱动。
\end{knowledgebox}

\subsection{实验：Planning 与 Execution 分开计分}

架构图说明了 SayCan 如何联合 usefulness 与 affordance，但还需要区分“计划在语言上合理”与“机器人最终完成任务”这两个不同命题。本节用独立指标检查高层排序和低层执行各自贡献；读结果时应先比较 planning 与 execution 的差距，再看移除 grounding 的 ablation。实验能够证明环境可执行性评分不可缺少，却不能单凭总体成功率判断剩余失败究竟来自感知、控制还是 skill coverage。

课堂报告 101 条 long-horizon instructions，规划成功率约 84\%，执行成功率约 74\%，并要求连续组合十个以上 navigation / manipulation skills。去掉 grounding 后性能接近减半，说明语言 plausibility 不能替代环境可执行性。

\lecturefigure{slide-27-saycan-experiment-results.jpg}{SayCan experiment overview：84\% planning、74\% execution 与 grounding ablation}{官方录像恢复幻灯片 V027；Ahn et al., 2022}

\begin{warningbox}{Planning Success 不等于 Task Success}
一个步骤序列可以语义正确，却因感知、抓取或导航失败而无法执行。分开报告 planning 与 execution 能定位瓶颈；若只报最终任务率，就无法判断失败来自 LLM、affordance model 还是 low-level skill。
\end{warningbox}

\subsection{PaLM-SayCan：外部模型进步如何传导}

slide 28 横轴是 language model size，纵轴是 planning performance。PaLM-SayCan 随模型规模上升而改善，说明在机器人数据不变时，外部 LLM 进步可以通过 planner 接口传导到系统。蓝色 FLAN-SayCan 点也提醒 instruction tuning 与模型家族会混入 scale 对比。

\lecturefigure{slide-28-palm-saycan-scaling.jpg}{PaLM-SayCan：更强 LLM 提升高层规划表现}{官方录像恢复幻灯片 V028；Stanford CS25 Lecture 15}

更好的 prompt 和 chain-of-thought 还能处理“找一种不是苹果的水果”或多步骤清理请求。右侧示例把自然语言解释转成可调用技能序列，展示 LLM 的语义 decomposition；但每个序列仍要经过 affordance score 才能进入机器人。

\lecturefigure{slide-29-palm-saycan-prompting.jpg}{PaLM-SayCan prompting：CoT 与更强提示生成多步技能序列}{官方录像恢复幻灯片 V029；Stanford CS25 Lecture 15}

\teachervoice{讲者最看重 SayCan 的“leverage”：不重新采集机器人轨迹，只替换更好的 LLM 或 prompt，规划能力就能上升。但他同样强调，这种 leverage 只作用于 skill selection，不能补上 low-level controller 根本不会的动作。}

\subsection{本章小结}

SayCan 用 LM usefulness 与 affordance probability 的乘积连接语言先验和物理执行。它证明外部 LLM 进步能传导到机器人 planning，也把系统上限暴露得很清楚：grounded skill library、affordance model 和低层执行仍是不可跳过的瓶颈。

